% ===========================================================================
\section{Discussion}
\label{sec:discussion}
% ===========================================================================

\subsection{Implications for System Design}

Our results suggest concrete recommendations for distributed system designers:

\para{Configurable heartbeat intervals} Systems should expose heartbeat
interval and election timeout as first-class tuning parameters with clear
documentation of the resilience/detection trade-off. Currently, most systems
(etcd, TiKV, CockroachDB) compile these values into binary defaults.

\para{Adaptive failure detection} Systems using fixed-interval heartbeats
(Raft-based) show the narrowest danger zones. Moving to adaptive approaches
($\phi$-accrual, sliding-window variance tracking) provides 2--5$\times$
wider jitter tolerance with comparable detection latency.

\subsection{CI Integration}
\label{sec:ci-integration}

\system{}'s minimal recipes translate directly to CI regression tests:

\begin{lstlisting}[language={},caption={CI configuration using \system{} recipe.},label={lst:ci}]
# .github/workflows/resilience.yml
name: Resilience Boundary Check
on:
  schedule:
    - cron: '0 2 * * 1'  # Weekly Monday 2AM
  pull_request:
    paths: ['src/raft/**']  # On consensus changes

jobs:
  check-boundaries:
    runs-on: self-hosted
    steps:
      - uses: actions/checkout@v4
      - name: Verify etcd election boundary
        run: |
          faultforge minimize \
            --oracle oracles/etcd-election.yaml \
            --system specs/etcd.yaml \
            --initial-severity slow-5s \
            --assert-boundary-above 15ms
\end{lstlisting}

The \texttt{--assert-boundary-above} flag causes the CI job to \emph{fail}
if the converged boundary drops below the specified threshold, enabling
regression detection, configuration validation, and version comparison.

\subsection{Boundary Interpretation}
\label{sec:interpretation}

The converged severity from binary search is an \emph{upper bound} on the
true danger zone boundary, not the exact boundary itself.

\para{For highly sensitive systems} (etcd, ZooKeeper, MongoDB, TiKV,
CockroachDB): The true boundary is \emph{below} 19.5\,ms (confirmed by
\xinda{}'s grid showing reproduction at 100\,$\mu$s). Our 19.5\,ms result
means: ``with 8 binary search steps from 5000\,ms, we cannot distinguish
the boundary from zero; it is below our precision limit.'' Increasing
$B_{\text{mag}}$ to 16 would resolve to 0.076\,ms (76\,$\mu$s), consistent
with \xinda{}'s findings.

\para{For systems near the precision limit} (Cassandra, HBase, Kafka,
Hadoop): The converged value approximates the true boundary within
$\pm S_{\max}/2^{B_{\text{mag}}}$. Redis (2.7\,s from 10\,s range,
precision 39\,ms) provides a tight bound.

\para{Practical significance} For operational purposes, the precision limit
(19.5\,ms) is sufficient to conclude ``this system is vulnerable to
negligible network perturbation.'' The exact boundary (whether 5\,$\mu$s or
50\,$\mu$s) does not change the operational response: both are below
any reasonable monitoring or alerting threshold.

\subsection{Limitations}
\label{sec:limitations}

\para{Single-node injection} Our evaluation primarily targets single-node
fault injection. Production environments may have correlated multi-node
degradations (rack-level network issues, storage array slowdowns).
Our multi-fault experiments (Section~\ref{sec:multi-fault-eval}) partially
address this, but full correlated-failure modeling remains future work.

\para{Containerized deployment} All systems run in Docker containers on a
single host. This eliminates real network hardware effects (switch buffers,
NIC queuing) and may under-represent multi-host timing variations. However,
container deployment is the standard for development testing and CI, which is
our primary target use case.

\para{Workload sensitivity} We use standard benchmarks (YCSB where applicable,
system-specific otherwise). Production workloads may have different access
patterns, sizes, and concurrency levels that shift fault boundaries.

\para{Monotonicity violations} While we validate monotonicity empirically
(Section~\ref{sec:monotonicity}), pathological workloads could violate it.
Example: a system with timeout-based retries might \emph{succeed} at very high
delay (because the retry timer has ample time to fire) but \emph{fail} at
moderate delay. We did not observe this in our evaluation but acknowledge it as
a theoretical limitation.

\para{Oracle sensitivity} Our results depend on oracle quality. The oracle
evaluation (Section~\ref{sec:oracle-eval}) quantifies this risk: 97\% recall
means approximately 3\% of actual failures go undetected.

\subsection{Threats to Validity}
\label{sec:threats}

\para{Internal validity} Non-determinism in distributed systems means
trial outcomes at boundary severities may vary. We mitigate this with:
(a)~repeated runs showing $<$5\% CV (Section~\ref{sec:stability}),
(b)~dense monotonicity probes (Section~\ref{sec:monotonicity}), and
(c)~the minimizer's conservative approach (using $hi$ as the converged
value, which always reproduces).

\para{External validity} Our 10 systems span multiple architectures (Raft,
ZAB, gossip, $\phi$-accrual, coordinated sessions) and languages (Go, Java,
C, Rust), but may not represent all production systems.

\para{Construct validity} We define ``danger zone boundary'' as the minimal
severity that triggers oracle-detectable symptoms. This may differ from the
minimal severity that triggers \emph{user-visible} impact (which may be higher
due to client-side retries, load balancer failover, etc.).

\subsection{Data Availability}

\system{} is released as open source. All oracle specifications, system
configurations, and experiment scripts are included in the repository.
Raw trial logs and converged boundaries for all 50 system--fault
configurations are available in \texttt{experiments/RESULTS.md}. We
encourage artifact evaluation and independent reproduction.
